% Generated draft: requires mathematical and source review.
\documentclass[11pt]{article}
\usepackage[margin=25mm]{geometry}
\usepackage{fontspec}
\setmainfont{Latin Modern Roman}
\usepackage{amsmath,amssymb,mathtools,unicode-math}
\setmathfont{Latin Modern Math}
\usepackage{enumitem,needspace,xurl,hyperref}
\setlength{\parindent}{0pt}
\setlength{\parskip}{5pt}
\setlength{\emergencystretch}{3em}
\newcommand{\N}{\mathbb{N}}
\newcommand{\Z}{\mathbb{Z}}
\newcommand{\Q}{\mathbb{Q}}
\newcommand{\R}{\mathbb{R}}
\newcommand{\C}{\mathbb{C}}
\newcommand{\F}{\mathbb{F}}
\newcommand{\A}{\mathbb{A}}
\newcommand{\PP}{\mathbb{P}}
\newcommand{\E}{\mathbb{E}}
\newcommand{\eps}{\varepsilon}
\newcommand{\dd}{\,\mathrm{d}}
\begin{document}
\section*{UnsolvedMath ID: 800631}
\textbf{Source problem code:} catalogue entry
\subsection{Definitions and mathematical statement}
The supplied statement has two unresolved prerequisites: its displayed ambient set \(\mathbb N\) cannot literally contain the strictly increasing functions under discussion, and neither \(Q\) nor the topology on \(H=B\cup Q\) is defined. The following gives a conditional formulation using the interpretation that the elements of \(B\) belong to \(\mathbb N^{\mathbb N}\). This interpretation requires confirmation, and no definition of \(Q\) is inferred.

For \(f,g\in\mathbb N^{\mathbb N}\), define eventual domination by
\[
 f\leq^*g\quad\Longleftrightarrow\quad
 \{n\in\mathbb N:f(n)>g(n)\}\text{ is finite}.
\]
A family \(A\subseteq\mathbb N^{\mathbb N}\) is unbounded if
\[
 \forall g\in\mathbb N^{\mathbb N}\ \exists f\in A\quad f\not\leq^*g.
\]
The bounding number is
\[
 \mathfrak b=\min\{|A|:A\subseteq\mathbb N^{\mathbb N}\text{ is unbounded}\}.
\]
A \(\mathfrak b\)-scale as described in the supplied record is an unbounded family
\[
 B=\{f_\alpha:\alpha<\mathfrak b\}\subseteq\mathbb N^{\mathbb N}
\]
of strictly increasing functions, meaning \(f_\alpha(n)<f_\alpha(n+1)\) for every \(n\in\mathbb N\), such that
\[
 \alpha<\beta<\mathfrak b\quad\Longrightarrow\quad f_\alpha\leq^*f_\beta.
\]

For a topological space \(X\), let \(\mathcal O(X)\) be the family of open covers of \(X\), with the convention that a cover does not contain \(X\) itself. An open \(\gamma\)-cover is an infinite family \(\mathcal U\) of open subsets of \(X\) such that
\[
 \forall x\in X\quad\{U\in\mathcal U:x\notin U\}\text{ is finite}.
\]
Write \(\Gamma(X)\) for the family of these covers. The property \(\mathsf S_1(\Gamma,\Gamma)\) means
\[
 \forall (\mathcal U_n)_{n\in\mathbb N}\in\Gamma(X)^{\mathbb N}\quad
 \exists (U_n)_{n\in\mathbb N}\quad
 \bigl[(\forall n\in\mathbb N\; U_n\in\mathcal U_n)
 \ \land\ \{U_n:n\in\mathbb N\}\in\Gamma(X)\bigr].
\]
In particular, the selected family must be infinite as a set of distinct open subsets.

For the contextual assertion, \(\mathsf U_{\mathrm{fin}}(\mathcal O,\Gamma)\) means that for every sequence \((\mathcal U_n)_{n\in\mathbb N}\) of open covers, each having no finite subcover, there are finite families \(\mathcal F_n\subseteq\mathcal U_n\) such that
\[
 \left\{\bigcup\mathcal F_n:n\in\mathbb N\right\}\in\Gamma(X).
\]

\textbf{Question, conditional on completing the source definitions.} For the \(\mathfrak b\)-scale \(B\) described above, form \(H=B\cup Q\), with \(Q\) and the topology on \(H\) to be recovered from the cited source. Must \(H\) satisfy \(\mathsf S_1(\Gamma,\Gamma)\)? Equivalently, must every sequence of open \(\gamma\)-covers of that space admit one selection from each cover whose selected family is again an open \(\gamma\)-cover?

The supplied record attributes to reference [ideals] the assertion that \(H^k\) satisfies \(\mathsf U_{\mathrm{fin}}(\mathcal O,\Gamma)\) for every positive integer \(k\), where finite powers carry the product topology. This is retained as supplied background, subject to the missing definitions and hypotheses; it is not independently verified here. The record also leaves unclear whether the question concerns every such scale or a particular scale fixed in preceding source context.
\subsection{Short English statement}
The question asks whether the space \(H=B\cup Q\), constructed from a \(\mathfrak b\)-scale, has the property that one can choose a single member of each open \(\gamma\)-cover in any sequence so that the chosen family is again a \(\gamma\)-cover. A precise instance requires the missing definition of \(Q\), the topology, and confirmation of the intended function space and scale quantifier. The supplied background attributes a weaker finite-selection property to all finite powers of \(H\).
\subsection{Sources}
\begin{itemize}
\item Local UnsolvedMath registry, permanent ID \(800631\), exact source problem code \texttt{catalogue entry}; supplied title: ``Selection Principles Milestone Problems --- Problem 3 ([ideals, Problem 17], [o-bdd, Problem 2]).'' The registry attributes the problem to Tsaban, with proposed year 2003, and labels its status as open. These are supplied registry labels, not independently verified conclusions.
\item Supplied source list: Tsaban, \emph{Selection Principles in Mathematics: A Milestone of Open Problems} (2003), Problem 3. Supplied URL: \url{https://arxiv.org/abs/math/0312182}. The record reports access on 2026-10-16 and extraction by \texttt{arxiv\_latex}; these provenance details are supplied and do not represent source access performed for this draft.
\item Supplied attribution/context: [ideals, Problem 17] and [o-bdd, Problem 2]. The supplied statement attributes the finite-power \(\mathsf U_{\mathrm{fin}}(\mathcal O,\Gamma)\) result to [ideals]. Full bibliographic entries for these reference labels are absent from the supplied record.
\end{itemize}
\end{document}
